% p1-appA-reproducibility

\section{P1 Appendix A: Reproducibility Protocol Checklist}

Appendix A: Reproducibility Protocol Checklist

        Figure (fig-p1-appA-reproducibility). Reproducibility triptych --- Checklist / Dataset cabinet /
                                                Enumerator.

  A.1 Dataset Preparation
  - [ ] P1: Constants T drawn exclusively from PDG 2024 and/or CODATA 2022, edition
  fixed before search begins.

  - [ ] P2: All constants in T are dimensionless.
  - [ ] P3: Only constants with relative experimental uncertainty better than 10-2 are
  included in the primary analysis.
  - [ ] P4: The dataset list (with values and uncertainties) is committed to version control
  before any symbolic search is executed, with a cryptographic hash.
  - [ ] P5: A separate out-of-sample set Toos is defined and frozen prior to analysis.
  - [ ] P6: All data files, version hashes, and Zenodo provenance records are deposited as
  specified in PhD: App.G Data Availability.

  A.2 Grammar Enumeration
  - [ ] G1: Grammar G$\varphi$ is implemented as in equations (1)--(2), with |theta|1 = |k| + |p| +
  |m| + |q|.
  - [ ] G2: Exponent bounds satisfy |k|, |p|, |m|, |q| $\leq$ Cmax with Cmax $\leq$ 15.
  - [ ] G3: The full expression set S(C) is deterministically enumerated (not sampled) for
  all C $\leq$ 10.
  - [ ] G4: The cardinality |S(C)| is explicitly reported for each complexity level.
  - [ ] G5: Numerical evaluation of F(theta) uses IEEE 754 double precision or higher.

  A.3 Residual Computation
  - [ ] R1: Approximation residual is defined as dC(X) = minF $\in$ S(C) |F - X| / |X|.
  - [ ] R2: Best-fit expression F\textasciicircum{}*(X) is recorded alongside the residual for all constants.
  - [ ] R3: All residuals are evaluated before any comparison against null models or
  selection of reported results.

  A.4 Null Model Evaluation
  - [ ] N1: Monte Carlo sampling for M0 uses at least B $\geq$ 10{,}000 independent samples.
  - [ ] N2: Permutation sampling for Mperm uses at least B $\geq$ 10{,}000 random
  permutations.
  - [ ] N3: Randomized grammar Mrand uses exponent vectors sampled uniformly from
  {-C, …, C}4.
  - [ ] N4: Null distributions are reported as quantiles (5th, 50th, 95th percentile) at each
  (C, $\varepsilon$) grid point.
  - [ ] N5: All 50 negative control realizations are evaluated and reported.

  A.5 Multiple Testing
  - [ ] T1: Bonferroni-corrected p-values computed using Neff(C) = |S(C)| as trial count.
  - [ ] T2: Family-wise tests use N $\cdot$ |S(C)| as the trial count.
  - [ ] T3: BH-adjusted p-values reported for the full dataset.
  - [ ] T4: No result reported as statistically significant unless it survives both analytic
  Bonferroni correction and permutation calibration.

  A.6 MDL Computation
  - [ ] M1: Model description length L(M1) computed as in equations (22)--(24).
  - [ ] M2: Results reported for at least three values of beta $\in$ {10, 100, 1000}.
  - [ ] M3: DeltaMDL reported with uncertainty from finite Monte Carlo sample.
  - [ ] M4: Positive compression claim requires DeltaMDL > 0 for all beta values in the
  range.

  A.7 Pellis Expansion
  - [ ] PE1: Expansion coefficients ck determined by iterative residual fitting as defined in
  Section 9.
  - [ ] PE2: Truncation depth K reported alongside expansion residual $\varepsilon$K(X) for each
  constant.
  - [ ] PE3: MDL cost of the Pellis expansion included via equation (40); only expansions
  with net DeltaMDL > 0 reported as compressive.

  A.8 Falsification Tests
  - [ ] F1: All five falsification criteria (Section 10.4) evaluated and reported.
  - [ ] F2: Out-of-sample test results computed and reported independently of primary
  analysis.
  - [ ] F3: Updated experimental values in out-of-sample tests cited with their
  PDG/CODATA reference year.
  - [ ] F4: All 50 negative control results reported in tabulated form.

